% Обозначение тиристора с управляющим электродом — CircuitTikZ, стиль л7
\documentclass[border=4pt]{standalone}
\usepackage{amsmath}
\usepackage{fontspec}
\setmainfont{Calibri}
\usepackage[european, RPvoltages]{circuitikz}
\usetikzlibrary{arrows.meta}
\definecolor{ink}{RGB}{31,42,51}
\definecolor{grp}{RGB}{46,110,140}
\ctikzset{bipoles/length=1.4cm}
\ctikzset{bipoles/thickness=0.9}
\ctikzset{diodes/fill=ink}
\ctikzset{inductor=american}

\begin{document}
\begin{tikzpicture}[line width=1.0pt, color=ink, font=\small]
  % --- тиристор: анод сверху, катод снизу (один путь — без швов)
  \draw (0,2.6) node[above]{$А$ (анод)}
        to[Ty, name=VS] (0,-0.4) node[below]{$К$ (катод)};
  % --- управляющий электрод (подпись у конца штриха затвора)
  \node[anchor=west, text=grp] at (VS.G) {УЭ};
  % --- ток
  \draw[-{Stealth}, thin, grp] (-0.6,1.7) -- (-0.6,0.8) node[midway, left, text=grp]{$i_{\mathrm{VS}}$};
\end{tikzpicture}
\end{document}
